\section{Discussion}\label{discussion}

\subsection{Quantitative and Ablation Interpretation}

The evaluation protocol separates seen-domain testing on Kvasir-SEG and
CVC-ClinicDB from external testing on CVC-300, CVC-ColonDB, and
ETIS-LaribPolypDB. All five test sets are isolated from checkpoint selection, which
uses only the internal validation subset of the combined development pool.
This design evaluates the sample-dependent guidance while retaining a convolutional
multi-scale pathway for efficient decoding.

FAFEM establishes the bottleneck feature-enhancement reference, MSCB-lite adds
multi-scale refinement at Stage 3, and RFG-MSCB adds adaptive branch modulation to
that path. Averaged across all five datasets, the final increment over FAFEM + MSCB
increases mDice from 0.869 to 0.887 and mIoU from 0.803 to 0.832 while reducing
MAE from 0.018 to 0.014. This improvement is consistent with a localized refinement
of the decoder. In the paired single-seed comparison, direct frequency guidance
without residual coupling produced a lower result than MSCB-lite.

The bounded residual design preserves the conventional MSCB response while allowing
sample-dependent branch modulation. Because residualization and frequency conditioning
were not evaluated independently, their individual contributions remain unresolved.
This adaptive refinement reuses the bottleneck descriptor at one decoder stage without
discarding the established multi-scale response or modifying the full decoder.

\subsection{Cross-Dataset Generalization}

CVC-300, CVC-ColonDB, and ETIS-LaribPolypDB are reserved for external
generalization evaluation and do not influence model development. Differences in
acquisition, annotation, case mix, and polyp morphology make these datasets a distinct
test of transfer beyond the Kvasir-SEG/CVC-ClinicDB development domain.
On the external datasets, the model obtained mDice/mIoU of 0.916/0.863 on
CVC-300, 0.824/0.747 on CVC-ColonDB, and 0.816/0.763 on ETIS-LaribPolypDB.
CVC-300 therefore had the strongest external result among the three evaluated
datasets and the highest displayed mDice and mIoU in its comparison table. On
CVC-ColonDB and ETIS, the model's mIoU was second-highest among the displayed
methods, whereas its mDice was below the two highest displayed values. These results
show that performance varied across the external datasets rather than supporting
uniform superiority. In the paired seed-42 ablation, the change from FAFEM + MSCB to
Residual RFG-MSCB likewise differed by dataset. Differences in acquisition,
annotation, and lesion appearance are plausible contributors, but were not isolated
experimentally. Stage-3 placement was selected on the development-domain evidence
and should not be assumed optimal for every external distribution.

Although the main comparison is restricted to methods with closely matched dataset
partitioning, differences in optimization, augmentation, backbone initialization,
model-selection criteria, and implementation details remain.

\subsection{Practical Trade-offs and Placement}

As summarized in Table~\ref{tab:model-complexity}, the final architecture has
30.06 M parameters and requires 27.82 GFLOPs at $352\times352$. Matched-hardware
latency was not measured.

The final model uses the new guidance at the Stage-3 decoder--encoder fusion. The
all-skip experiment was exploratory, single-seed, incomplete, and used a different
batch size. Deployment-oriented placement and latency optimization remain future work.

\subsection{Limitations}

\begin{itemize}
\item \textbf{No statistical significance testing:} All reported differences are descriptive; no hypothesis testing was performed.
\item \textbf{Software records:} Runs used an NVIDIA GeForce RTX 4090 and PyTorch 2.5.1+cu124 on Windows, but the complete dependency environment was not retained.
% TODO: replace the seed-wise training-duration limitation after the final three-seed run.
\item \textbf{Training duration:} With zero-based epoch indices, seed 7777 stopped after epoch 167 (168 completed epochs; patience 30), while seeds 42 and 6543 completed epochs 0--199; this asymmetry may affect the mean estimate.
\end{itemize}

\subsection{Future Work}

\begin{itemize}
\item Systematic multi-seed placement ablation with controlled batch size and learning rate.
\item Disentanglement of the RFG-MSCB components (residual form, initialization schedule, frequency conditioning) via factorial experiments.
\item Investigation of initialization sensitivity using completed default-initialization runs.
\item Domain adaptation or targeted augmentation to improve generalization on CVC-ColonDB and ETIS-LaribPolypDB.
\item Evaluate between-configuration differences using larger seed ensembles or bootstrap confidence intervals.
\item Record hardware and software configurations systematically in future experiments.
\end{itemize}
